\chapter{Finite-place fields and arithmetic comparison}
\label{v4-arith:drinfeld-kapranov}

A formal Heisenberg algebra exists over every characteristic-zero
field, including \(\mathbb Q_p\). Its insertion coordinate belongs
to an auxiliary curve over that field. A comparison with an
arithmetic space must specify the support, coefficient maps and
topology that relate them.

There are three separate choices. A logarithmic energy assigns
real weights to states. A coefficient construction determines the
ring over which local fields are assembled. A geometric realization
specifies the diagonal and its normal directions. We examine these
choices before stating the arithmetic comparison problem.

\section{The formal algebra over a finite-place field}
\label{v4-arith:dk:algebraic-step}

A vertex algebra over \(k\) has a state-field map
\[
                         Y:V\otimes_kV\longrightarrow V((z)),
\]
with a vacuum, translation and the formal field identities.
No absolute value occurs in those identities. The construction
in Section~\ref{v4-sec:singular-heisenberg} gives two orthogonal
level-one currents over \(\mathbb Q\), and scalar extension gives
\[
 \mathcal H_k=\operatorname{Sym}_k(k[t]a\oplus k[t]b),
 \qquad [a_m,a_n]=[b_m,b_n]=m\delta_{m+n,0},\quad[a_m,b_n]=0.
\]
The monomial \(t^{m-1}a\) denotes a mode-\(m\) creation state.

\begin{proposition}\label{v4-arith:prop:formal-padic-va}
For every characteristic-zero field \(k\), the displayed Fock
space is a rank-two formal Heisenberg vertex algebra, with
\[
 a(z)a(w)\sim(z-w)^{-2},\qquad
 b(z)b(w)\sim(z-w)^{-2},\qquad a(z)b(w)\sim0.
\]
In particular it gives a formal vertex algebra over \(\mathbb Q_p\).
\end{proposition}

\begin{proof}
Define negative modes by multiplication and positive modes by the
corresponding weighted derivatives. Their commutators give the
three contractions. Normal ordering and the finite Wick expansion
prove locality, translation and the vacuum identities, as in
Proposition~\ref{v4-prop:singular-heisenberg-fields}.
All formulas commute with extension from \(\mathbb Q\) to \(k\).
\end{proof}

The variable \(z\) remains formal. A norm on \(V\), a completion
of its tensors and a domain for \(z\) are additional analytic data.
Neither the existence nor failure of one analytic completion
changes this algebraic construction.

\section{Conformal normalization and logarithmic energy}
\label{v4-arith:dk:conformal-vector-scaling}

The standard conformal state is
\[
              \omega=\tfrac12(a_{-1}^2+b_{-1}^2)\mathbf1.
\]
Its field \(T\) has central charge two and integer oscillator
weights. For \(T'=\lambda T\), its second pole is
\(2\lambda^2T\), whereas the Virasoro normalization requires
\(2T'=2\lambda T\).

\begin{observation}\label{v4-arith:obs:log-p-not-padic}
On the nontrivial Heisenberg algebra, a scalar multiple
\(\lambda\omega\) implements the original conformal structure
only when \(\lambda=1\). The real weight \(\log p\) in
\(H_p=(\log p)L_0\) is external energy data.
It does not define a rescaled conformal state over \(\mathbb Q_p\).
\end{observation}

Indeed the second pole forces \(\lambda^2=\lambda\).
The value zero cannot give the required translation operator.
This leaves \(\lambda=1\). The fourth pole alone would give
\(2\lambda^2\), but it cannot determine a conformal structure
without the other poles.

The \(p\)-adic logarithm provides different data. Its convergent
series on principal units extends to \(\mathbb Q_p^\times\)
after one chooses its value on \(p\). The usual normalization
sets that value to zero. This convention is not an interpretation
of the real number \(\log p\), nor is it the only possible
choice of a logarithm on the valuation factor.

\begin{convention}\label{v4-arith:conv:log-p-archimedean}
Use the standard rank-two conformal state at each finite place.
When the coefficient space is complex and prime-labeled, use
\((\log p)L_0\) as a separately specified energy operator.
A comparison with a \(p\)-adic coefficient system must supply
its own coefficient maps and normalizations.
\end{convention}

\section{Common coefficients and restricted tensors}
\label{v4-arith:dk:restricted-tensor-not-field}

For finite sets of places, tensor products require a common
base ring and actual maps from that base. Since every
\(\mathbb Q_p\) is a \(\mathbb Q\)-algebra, the algebraic tensors
\[
                         R_S=\bigotimes_{p\in S,\mathbb Q}\mathbb Q_p
\]
are defined. Their transition maps insert units. They are
injective as maps of rational vector spaces.

\begin{observation}\label{v4-arith:obs:no-single-field}
There is an abstract common field extension of all the
\(\mathbb Q_p\). This fact does not provide canonical or
continuous coefficient embeddings, an adelic topology, or a
restricted tensor of representations.
\end{observation}

To construct such a field, take a maximal-ideal quotient of
\(\varinjlim_SR_S\), as in
Proposition~\ref{v4-prop:singular-common-coefficients}.
The transition injections make this a nonzero ring.
Each local field injects into the quotient because its map is
unital. This construction retains its algebraic scope.

Three coefficient constructions illustrate the further choices.
First, choose embeddings into a common characteristic-zero field
and extend the formal Heisenberg factors there. Even an abstract
embedding \(\mathbb Q_p\to\mathbb C\) is discontinuous, since it
sends the sequence \(p^r\to0\) to an unbounded real sequence.
Extending scalars in this way does not transport the local norm.

Second, work over an adelic coefficient ring with specified
modules and lattices. The canonical map is the projection
\(\mathbb A_{\mathbb Q}\to\mathbb Q_p\). It does not supply a
unital map in the reverse direction for scalar extension.
The finite-place inclusion supported only at \(p\) sends the
unit to a proper idempotent. Its module convention must therefore
be stated explicitly.

Third, use the rational model itself. The assignment
\[
                  k\longmapsto\mathcal H_{\mathbb Q}\otimes_{\mathbb Q}k
\]
is already a functor on characteristic-zero fields.
A field embedding acts on coefficients and preserves the field
operations. A descent datum for another form, or an arithmetic
Galois action with Frobenius compatibility, is additional input.

\begin{warning}
None of these coefficient constructions produces a Hecke action
or identifies the oscillator algebra with an arithmetic trace
class. Those comparisons require independently defined targets
and maps preserving their operations.
\end{warning}

\section{Completion in the prime and completion at the diagonal}
\label{v4-arith:dk:padic-diagonal-no-thickening}

Put \(R=\mathbb Z_p\otimes_{\mathbb Z}\mathbb Z_p\).
The \(p\)-adic completion of this ring and the formal completion
at its multiplication ideal use different ideals.
The first is computed by reduction modulo powers of \(p\).

\begin{proposition}\label{v4-arith:prop:padic-diagonal-collapse}
Multiplication induces
\[
        \varprojlim_r R/p^rR\cong\mathbb Z_p.
\]
The ordinary multiplication map \(R\to\mathbb Z_p\) has a
nonzero kernel. After inverting \(p\), its diagonal conormal
module is \(\Omega_{\mathbb Q_p/\mathbb Q}\ne0\).
\end{proposition}

\begin{proof}
At every finite precision,
\[
 R/p^rR\cong
     (\mathbb Z_p/p^r)\otimes_{\mathbb Z}(\mathbb Z_p/p^r)
                     \cong\mathbb Z/p^r.
\]
Taking this inverse limit proves the first assertion.
For \(a\in\mathbb Z_p\setminus\mathbb Q\), the element
\(a\otimes1-1\otimes a\) is nonzero before completion.
A rational linear functional taking \(1\) to zero and \(a\)
to one detects it after inverting \(p\).
Multiplication and every finite reduction kill it.

The conormal module of multiplication on
\(\mathbb Q_p\otimes_{\mathbb Q}\mathbb Q_p\) is the module
of Kähler differentials. A derivation along an element of a
transcendence basis is nonzero. Characteristic zero extends
that derivation across the remaining algebraic extension.
Thus this conormal module is nonzero, as in
Proposition~\ref{v4-prop:singular-padic-tensor}.
\end{proof}

Flatness does not make an ordinary tensor commute with this
inverse limit. The nonzero kernel just exhibited shows why
that interchange would give the wrong ordinary ring.

\begin{corollary}\label{v4-arith:cor:no-padic-ope}
The \(p\)-adic completion calculation does not imply that every
finite-place field operation is locally constant. On the auxiliary
line \(\mathbb A^1_{\mathbb Q_p}\), the rank-two currents have
nonzero principal-parts operations at the diagonal.
\end{corollary}

\begin{proof}
That diagonal has parameter \(u=z-w\), with formal normal
ring \(\mathbb Q_p[z][[u]]\). The current product has principal
part \(u^{-2}du\). This is nonzero in its full normal module.
Completion of the coefficient ring in \(p\) does not remove
this independent normal variable.
\end{proof}

\section{Choosing a local coefficient model}
\label{v4-arith:dk:E0-local-model}

A locally constant coefficient system is a possible model.
It must be specified rather than deduced from the preceding
completion formula. A commutative coefficient algebra with a
Hecke endomorphism is also different from the formal Heisenberg
algebra, whose current first product is nonzero.

\begin{definition}\label{v4-arith:def:E0-local-factor}
A constant local coefficient datum consists of a commutative
unital coefficient algebra \(B_p\), a chosen filtration or
topology, and a specified compatible endomorphism \(T_p\).
Its constant vertex operation is \(Y(b,z)c=bc\), with zero
translation. This is one possible local datum for a comparison
with an independently defined arithmetic object.
\end{definition}

An \(E_0\)-algebra by itself is pointed data, without a binary
multiplication supplied by that notation. The commutative product
in the definition is an explicit additional structure.
Its endomorphism is not simultaneously a Hecke and Adams
operation unless the corresponding relations are proved.

\begin{theorem}[Complex auxiliary realization]
\label{v4-arith:thm:archimedean-is-holomorphic}
The scalar extension \(\mathcal H_{\mathbb C}\) gives a
Heisenberg field algebra on the auxiliary complex line, with
its nonzero diagonal principal-parts operation.
\end{theorem}

\begin{proof}
The formal field construction and normal residue operation are
those of Sections~\ref{v4-sec:singular-heisenberg}
and~\ref{v4-sec:singular-residue-coalgebra} over \(\mathbb C\).
Translation gives differential linearity, and the field identities
give the residue identities.
\end{proof}

Identifying this auxiliary realization with a specified
archimedean arithmetic object requires a comparison map.
The coefficient extension alone does not define that object.

\section{The adelic and automorphic comparison}
\label{v4-arith:dk:true-hidden-structure}

A restricted tensor representation starts with representations
\((\pi_v,V_v)\) of the local groups over a common field.
At almost every finite place it also specifies a vector fixed
by a chosen compact open subgroup. The tensor transition maps
insert these vectors. A restricted product of group elements
then acts because only finitely many factors move them.

This constructs an abstract representation on its stated tensor
carrier. Realizing it as automorphic functions requires further
global data. In particular, local oscillator fields do not
already specify the \(GL_2(\mathbb Q_v)\)-representations or
the required global realization.

\begin{question}\label{v4-arith:thm:adelic-automorphic-placement}
Choose arithmetic local coefficient objects, local group actions,
and unramified reference vectors over a specified common base.
Construct their restricted tensor and an intertwiner to an
independently defined automorphic recipient.
Can this intertwiner be compared with the local field operations,
including their collision maps, vacuum data and coefficient topology?
\end{question}

\begin{remark}
A compatible adelic system may be a useful geometric recipient
for this question. Its structure must determine the relevant
local and global maps. Neither the coefficient-field choices
nor the completion of \(\mathbb Z_p\otimes\mathbb Z_p\)
force that recipient or prove that singularities occur only
at the archimedean place.
\end{remark}

\section{Formal loop spaces and the arithmetic problem}
\label{v4-arith:dk:kapranov-built}

\begin{remark}
Kapranov and Vasserot construct formal loop spaces for schemes
of finite type over a field. Their global versions form
factorization monoids. For smooth schemes their geometric
construction recovers the chiral de Rham complex and its
vertex structure.\footnote{M.~Kapranov and E.~Vasserot,
\emph{Vertex algebras and the formal loop space},
Introduction, pp.~1--2,
\href{https://arxiv.org/abs/math/0107143}{arXiv:math/0107143}.}
The field-based construction supplies a comparison model.
An arithmetic extension must define its own loop objects,
coefficient categories and factorization maps.
\end{remark}

\begin{remark}
Reciprocity and cohomology determinants suggest further comparison
data. A determinant line requires a specified perfect complex,
and a reciprocity statement requires its actual local symbols
and common recipient. An analogy with Langlands correspondence
or a field theory does not provide these constructions.
The automorphic problem above remains one possible formulation.
\end{remark}

\section{The local and global maps still required}
\label{v4-arith:dk:summary}

The formal finite-place Heisenberg algebra, its auxiliary diagonal
operation, and the \(p\)-adically completed coefficient tensor
are explicit constructions. The ordinary tensor and diagonal
conormal module retain the information that the \(p\)-adic
completion forgets. Logarithmic energy preserves its separate
real normalization.

A geometric arithmetic realization must relate these specified
carriers through coefficient and support maps. An automorphic
realization must additionally supply its group actions and global
intertwiner. The choice of a common field, a locally constant
model, or an adelic tensor leaves those maps to be constructed.
